% Avsnitt 14.1-14.6, ur lectures/L14/appendix/a_tx_shift_reg.md.

\section{Gränssnitt}
\label{c14:sec:iface}

\modulefig[0.62]{lectures/L14/appendix/images/tx_shift_reg.png}

\begin{booktable}{@{}p{5.6em}p{3.4em}p{10em}L@{}}
  \thead{Port} & \thead{Riktn.} & \thead{Typ} & \thead{Betydelse} \\
  \midrule
  \code{clock}, \code{reset_s2_n} & in & \code{std_logic} & 50 MHz klocka, synkroniserad aktiv låg
    reset. \\
  \code{load} & in & \code{std_logic} & Låser \code{data}/\code{bit_count} som nästa grupp
    \textbf{och} presenterar dess första bit den här cykeln. \\
  \code{data} & in & \code{byte_t} & Gruppen, med sina \code{bit_count} giltiga bitar packade i de
    \textbf{mest} signifikanta bitarna (första biten som sänds är bit 7). \\
  \code{bit_count} & in & \code{std_logic_vector(3 downto 0)} & Antal giltiga bitar i \code{data},
    1--8. \\
  \code{shift} & in & \code{std_logic} & Pulsar en gång per CAN-bit (från
    \code{bit_timer.bit_done}) för att driva strömmen framåt. \\
  \code{tx_bit} & out & \code{std_logic} & Registrerad bit att driva under den kommande
    bitperioden. \\
  \code{stuff} & out & \code{std_logic} & \code{'1'} på just den skiftning som lägger ut en
    stoppbit; maskar bort den biten ur CRC-matningen. \\
  \code{done} & out & \code{std_logic} & Pulsar när gruppen är helt utsänd (se tajmingen). \\
  \code{bit_valid} & out & \code{std_logic} & \code{'1'} den enda cykel då en ny bit lades ut på
    \code{tx_bit}. \\
\end{booktable}

\code{tx_bit} är den enda utgången som är en \emph{nivå}, hållen tills nästa bit presenteras.
\code{stuff}, \code{done} och \code{bit_valid} är alla \textbf{pulser på en cykel}: ge dem förvalet
\code{'0'} högst upp i den klockade processen och höj dem bara i den gren som förtjänar dem, samma
idiom som \code{bit_timer} använder för \code{sample} och \code{bit_done}.

Det spelar störst roll för \code{bit_valid} och \code{stuff}, eftersom \code{can_controller} bygger
\code{crc15}s enable av paret, som \code{bit_valid and not stuff} (\chapref{c16:ch}):
\code{bit_valid} säger att en bit lades ut, och \code{stuff} maskar bort de utlagda bitar som CRC:n
aldrig får se, eftersom CRC:n bara omfattar riktiga bitar. Håll \code{bit_valid} som en nivå i
stället, så integrerar motorn samma bit vid varje klockflank under bitperioden, femtio gånger om, och
förstör CRC:n i tysthet.

\code{load} och \code{shift} utesluter varandra inom en cykel; \code{can_controller} aktiverar aldrig
båda. \code{tx_shift_reg_tb} binder \textbf{positionellt}, så ordningen och typerna ovan är det som
måste stämma; namnen är dina.

% ------------------------------------------------------------------------------------------
\section{Vad modulen måste minnas}
\label{c14:sec:state}

Börja där inledningen pekade: pappersversionen av stoppbitsregeln såg \emph{tillbaka över följden},
och hårdvaran har ingen följd att se tillbaka över. Så innan någon process skrivs, lista vad som
måste bo i ett register. Allt nedan följer av modulens två uppgifter \textemdash{} att hålla en grupp
i rörelse, och att hålla stoppbitsregeln ärlig \textemdash{} och varje regel i resten av det här
kapitlet talar i termer av exakt de här fem signalerna:

\begin{vhdlcode}
signal shift_reg    : byte_t;                     -- Remaining chunk bits, next one in bit 7.
signal bits_left    : natural range 0 to 8;       -- Real bits of the chunk still to present.
signal last_bit     : std_logic;                  -- The bit most recently placed on the wire.
signal consecutive  : natural range 0 to MAX_RUN; -- Current run length; 0 = none.
signal stuff_pending: std_logic;                  -- A stuff bit is owed on the next shift.
\end{vhdlcode}

\begin{itemize}
  \item \textbf{\code{shift_reg} och \code{bits_left} är gruppen i flykt.} Registret håller de
    osända bitarna vänsterjusterade så att ”nästa bit” alltid är \code{shift_reg(7)};
    \code{bits_left} är hur många riktiga bitar som återstår. Mer grupphållning än så finns inte.
  \item \textbf{\code{last_bit} och \code{consecutive} är stoppbitsregelns minne}, och tillsammans
    kallas de nedan för \textbf{följdspåraren}: biten som senast sändes, och hur många lika bitar i
    rad den avslutade. \code{MAX_RUN} lika bitar \textemdash{} fem, en konstant som \code{can_def}
    tillhandahåller, eftersom länkens båda ändar måste vara exakt överens om den \textemdash{} är en
    full följd; paret är allt regeln någonsin behövde av att ”se tillbaka”. Båda skrivs på exakt ett
    ställe, en procedur som processen deklarerar för ändamålet (nedan).
  \item \textbf{\code{stuff_pending} är regelns enbitsskuld.} När en följd når fem sänds stoppbiten
    \emph{inte} omedelbart \textemdash{} den femte biten behöver fortfarande sin egen fulla period på
    bussen. Att nå fem \emph{armerar} därför bara en insättning, och den här flaggan bär skulden
    vidare till nästa \code{shift}.
\end{itemize}

Utgångarna registreras i samma process: \code{tx_bit} är en nivå, hållen tills nästa bit ersätter
den; \code{stuff}, \code{done} och \code{bit_valid} är encykelspulserna som \secref{c14:sec:iface}
beskrev, som förval låga högst upp i processen och höjda bara i den gren som förtjänar dem.

% ------------------------------------------------------------------------------------------
\section{Processen, uppifrån och ned}
\label{c14:sec:process}

En klockad process håller allt ovanstående. Genomgången nedan följer koden i den ordning den skrivs,
så varje rad har ett hem innan nästa dyker upp; kodfragmenten är processens deklarativa del följd av
på varandra följande skivor av en enda \code{if}/\code{elsif}-kedja, och lagda i följd inuti
skelettet nedan, vars resetgren du fyller i utifrån dess enradiga beskrivning, är de hela processen.
En rad använder \code{to_integer(unsigned(...))}, så filen behöver \code{use ieee.numeric_std.all;}
vid sidan av \code{ieee.std_logic_1164} och \code{work.can_def}.

\textbf{Skelettet.} Processen deklarerar sin egen procedur före \code{begin}; sedan asynkron reset
först, allt annat på den stigande flanken:

\begin{vhdlcode}
process(clock, reset_s2_n) is
    -- One procedure, the run tracker (below).
begin
    if (reset_s2_n = '0') then
        -- Clear everything.
    elsif (rising_edge(clock)) then
        -- Defaults, then the load branch or the shift branch.
    end if;
end process;
\end{vhdlcode}

\subsection{Följdspåraren, skriven en gång}
\label{c14:sub:tracker}

Varje gren som lägger ut en bit behöver samma få rader, så de bor i en procedur som deklareras i
processens egen deklarativa del, mellan \code{process(...) is} och \code{begin}:

\begin{vhdlcode}
    -- The run tracker: runs on the bit just placed on the wire.
    procedure track_run(new_bit: in std_logic) is
    begin
        if ((new_bit = last_bit) and (consecutive /= 0)) then
            consecutive <= consecutive + 1;
            if (consecutive = MAX_RUN - 1) then -- Run reaches MAX_RUN.
                stuff_pending <= '1';
            end if;
        else
            consecutive <= 1;                   -- Start a new run.
        end if;
        last_bit <= new_bit;
    end procedure;
\end{vhdlcode}

Det här är bokens första underprogram, så fyra saker om det innan något anropar det:
\begin{itemize}
  \item \textbf{Den når signalerna direkt.} \code{last_bit}, \code{consecutive} och
    \code{stuff_pending} är arkitektursignaler, synliga inuti proceduren eftersom proceduren är
    deklarerad inuti processen, som ligger inuti arkitekturen. De är inte parametrar och ska inte
    vara det; bara den \emph{utlagda biten} skiljer sig mellan anroparna, och det är den enda
    parameter som finns.
  \item \textbf{Att deklarera den inuti processen är det som gör signaltilldelningarna lagliga
    alls.} Ett underprogram får tilldela en signal bara i två situationer: signalen är en formell
    parameter av klassen \code{signal}, eller underprogrammet är deklarerat inuti en process, och då
    använder dess tilldelningar den processens drivare, precis som om raderna hade skrivits på varje
    anropsställe. Den här proceduren vilar på det andra fallet. Flytta upp den oförändrad till
    arkitekturens deklarativa del så slutar den kompilera, ett fel per tilldelning:
    \code{signal "consecutive" is not a formal parameter}. Att träda \code{consecutive},
    \code{last_bit} och \code{stuff_pending} genom som formella parametrar av klassen \code{signal}
    skulle uppfylla regeln, till priset av tre extra argument vid varje anrop och en verklig fråga om
    vilken process som äger drivarna. Deklarerad här kan bara den här processen anropa den, så bara
    den här processen driver dem.
  \item \textbf{Tajmingen är oförändrad.} En \code{in}-parameter evalueras vid anropet, så
    \code{new_bit} bär det värde före flanken som anroparen skickade in, och signalläsningarna inuti
    kroppen är också från före flanken (regeln om variabler kontra signaler från \chapref{c5:ch}).
    Proceduren beter sig precis som samma rader skrivna inline.
  \item \textbf{Syntesen inlinar den.} Resultatet är samma logik replikerad under varje anropares
    grenvillkor; ingenting delas, multiplexas eller anropas vid körning. Det här är faktorisering på
    källkodsnivå, och nätlistan blir identisk hur som helst.
\end{itemize}

Två detaljer inuti kroppen, och en om var den anropas:
\begin{itemize}
  \item \textbf{\code{consecutive = MAX_RUN - 1} betyder ”följden är nu full”.} Läsningen ger värdet
    före flanken, så jämförelsen sker mot följdlängden \emph{innan} den här biten förlängde den. Den
    semantiken är precis rätt här, och det är därför koden jämför mot \code{MAX_RUN - 1} och inte
    \code{MAX_RUN}.
  \item \textbf{Håll reda på vilket \code{if} som äger vilken gren.} \code{consecutive <= 1} är det
    \emph{yttre} \code{if}-satsens \code{else} (varje utlagd bit som bryter följden startar en ny),
    och \code{last_bit <= new_bit} står efter hela \code{if}-satsen och uppdateras vid \emph{varje}
    utläggning, lika eller inte. Häng någon av dem på den inre \code{MAX_RUN - 1}-kontrollen, så
    startar en bruten följd aldrig om, eller så fryser \code{last_bit}.
  \item \textbf{Anropa den från en gren som lägger ut en bit, aldrig från processens toppnivå.}
    Ingenting på ledningen ändras under de dryga fyrtionio klockflanker som ligger mellan
    presentationerna, så ett anrop som körs vid varje flank räknar spökbitar. Grupper vet den
    ingenting om åt något håll: ledningen har ingen aning om var en grupp slutar och nästa börjar, så
    det har inte spåraren heller.
\end{itemize}

\subsection{Reset och förval}
\label{c14:sub:defaults}

\textbf{Reset nollställer allt}: utgångarna, gruppen (\code{shift_reg}, \code{bits_left}) och
följdspåraren (\code{last_bit}, \code{consecutive}, \code{stuff_pending}). Det här är det
\textbf{enda} stället spåraren någonsin nollställs. \code{load} får inte röra den: stoppningen ser
den utsända strömmen som sammanhängande över hela ramen, så en följd som avslutar en grupp och
fortsätter in i nästa tvingar fortfarande fram en stoppbit precis vid gränsen. (Det är scenariot med
kedjad omladdning, test 3 i testbänken: en grupp slutar med fyra \code{'1'}:or, nästa öppnar med en
\code{'1'}, och stoppbiten landar före den gruppens andra riktiga bit.)

\textbf{Förvalen härnäst, och bara de här tre.} Pulsutgångarna faller vid varje klockflanks början,
\code{bit_timer}-idiomet, så grenarna nedan bara \emph{höjer} dem:

\begin{vhdlcode}
stuff     <= '0';
done      <= '0';
bit_valid <= '0';
\end{vhdlcode}

\begin{warning}
\code{stuff_pending} hör \textbf{inte} hemma här, hur mycket den än liknar dem. Den är tillstånd,
inte en puls: armerad vid den skiftning som fullbordar en följd, betald vid \emph{nästa} skiftning,
en hel bitperiod (femtio klockflanker) senare. Ge den förvalet lågt här, så avdunstar skulden vid
den första lediga flanken däremellan, och då stoppar modulen helt enkelt aldrig in något. Den skrivs
på exakt tre ställen: nollställd vid reset, armerad inuti \code{track_run}, och nollställd av den
gren som betalar den.
\end{warning}

\subsection{Laddgrenen}
\label{c14:sub:load}

\textbf{\code{load = '1'} presenterar den första biten omedelbart.} En bitperiod börjar i samma
ögonblick som den föregående gruppens sista \code{bit_done} går, vilket också är när nästa grupps
\code{load} går. Om \code{load} bara låste data och väntade på den första \code{shift}-pulsen skulle
den första perioden fortfarande visa den föregående gruppens sista bit, en hel bit försent. Så grenen
lägger ut \code{data(7)} nu, för registret i förväg förbi den, sätter \code{bits_left} till
\code{bit_count - 1}, och kör sedan den utlagda biten genom följdspåraren:

\begin{vhdlcode}
if (load = '1') then
    tx_bit    <= data(7);                 -- Output the first bit immediately.
    shift_reg <= data(6 downto 0) & '0';  -- Remove the output bit and prepare the next one.
    bits_left <= to_integer(unsigned(bit_count)) - 1;  -- bit_count is 1-8 by contract:
                                                       -- 0 would underflow bits_left here.
    bit_valid <= '1';                     -- Indicate that a valid bit has been output.

    track_run(data(7));                   -- data(7) was just placed on the wire.
\end{vhdlcode}

Anropet är varken valfritt eller en detalj: \code{load} lägger ut en bit på ledningen, så spåraren
måste se den som vilken annan som helst. Utelämna det, så ligger följdräkningen en bit efter under
resten av ramen; test 1 i testbänken fångar det vid skiftning 5, där stoppbiten som de fem
\code{'1'}:orna förtjänat uteblir.

\subsection{Skiftgrenen}
\label{c14:sub:shift}

\textbf{\code{shift = '1'} gör exakt en av tre saker} (aldrig i samma cykel som \code{load}; det är
därför grenen är ett \code{elsif}), prövade i den här prioritetsordningen. Först den stoppbit som är
skyldig, om det finns någon:

\begin{vhdlcode}
elsif (shift = '1') then
    if (stuff_pending = '1') then
        tx_bit        <= not last_bit;  -- The stuff bit inverts the run.
        stuff         <= '1';
        bit_valid     <= '1';
        stuff_pending <= '0';           -- The debt is paid.

        track_run(not last_bit);        -- A stuffed bit is a placed bit like any other.
\end{vhdlcode}

En stoppbit går genom spåraren precis som en riktig, och att skicka in \code{not last_bit} är det som
gör att det blir rätt utan något specialfall: inuti kroppen är jämförelsen \code{new_bit = last_bit}
då falsk per konstruktion, så anropet tar \code{else}-vägen och gör precis det regeln vill, startar
om följden på 1 med stoppbitens eget värde i \code{last_bit}. Ingenting armerar någon ny skuld här
heller, eftersom armeringen sitter i den gren som inte kan tas. Gruppen själv gör inga framsteg:
\code{shift_reg} och \code{bits_left} lämnas orörda.

Härnäst den \emph{avslutande} skiftningen:

\begin{vhdlcode}
    elsif (bits_left = 0) then
        done <= '1';                    -- The closing shift changes nothing else.
\end{vhdlcode}

Varje bit har lagts ut och den sista har nu fått sin fulla period på bussen; \code{tx_bit} i
synnerhet behåller sitt värde. I annat fall: sänd nästa riktiga bit, samma form som
\code{load}-grenen med \code{shift_reg(7)} i \code{data(7)}:s ställe:

\begin{vhdlcode}
    else
        tx_bit    <= shift_reg(7);            -- Output the next bit.
        shift_reg <= shift_reg(6 downto 0) & '0';
        bits_left <= bits_left - 1;
        bit_valid <= '1';

        track_run(shift_reg(7));              -- May arm a stuff bit for the next shift.
    end if;
end if;
\end{vhdlcode}

Prioritetsordningen är bärande: en väntande stoppbit går före den avslutande skiftningen, och det är
det som får specialfallet nedan att fungera utan någon extra kod; en grupp vars \emph{sista} riktiga
bit fullbordar en följd får ändå sin stoppbit utsänd innan \code{done} går. (\code{load} har aldrig
någon utestående stoppbit att bekymra sig om: \code{can_controller} laddar bara efter \code{done},
och \code{done} kommer efter att varje väntande stoppbit har sänts.)

% ------------------------------------------------------------------------------------------
\section{Utgångarnas tajming, den del testbänken låser exakt}
\label{c14:sec:timing}

\begin{itemize}
  \item \code{tx_bit} sätts av \textbf{både} \code{load} och \code{shift}; den avslutande skiftningen
    lämnar den oförändrad.
  \item \code{stuff} pulsar bara vid den skiftning som faktiskt lägger ut en stoppbit, aldrig vid
    \code{load} (en grupps första bit kan \emph{utlösa} en stoppbit, men den stoppbiten sänds vid en
    senare skiftning).
  \item \code{done} pulsar \textbf{inte} vid den skiftning som presenterar gruppens sista riktiga bit
    eller stoppbit; den biten behöver först sin egen fulla period på bussen. Den pulsar en skiftning
    senare, vid den avslutande skiftningen, som inte presenterar något nytt. Att låta \code{done} gå
    tidigt skulle låta \code{can_controller} ladda nästa grupp en hel bitperiod för tidigt.
  \item \code{bit_valid} är \code{'1'} vid \code{load} och vid varje skiftning som lägger ut en ny
    bit (riktig eller stoppbit), och \code{'0'} vid den avslutande skiftningen. Konsumenter som för
    \code{tx_bit} vidare bit för bit grindar på \code{bit_valid}, inte på ”load eller shift”, annars
    skulle de sampla den avslutande skiftningens oförändrade \code{tx_bit} som en dubblett.
    \code{crc15} är den konsument som spelar roll, och den behöver \code{stuff} också:
    \chapref{c16:ch} grindar dess enable på \code{bit_valid and not stuff}, eftersom en stoppbit
    läggs ut på ledningen men aldrig får gå in i CRC:n.
  \item Specialfall: om \code{bit_count = 1} och den enda biten utlöser en stoppbit, väntar
    \code{done} ändå på att stoppbiten sänts först.
\end{itemize}

% ------------------------------------------------------------------------------------------
\section{En grupp, presentation för presentation}
\label{c14:sec:trace}

Fyra tajmingregler samspelar här, och prosa är ett dåligt sätt att se dem. Nedan följer en komplett
grupp: \code{data = "11111000"}, \code{bit_count = 8}, laddad in i ett nyss resettat register. Varje
rad är en \emph{presentation}, antingen \code{load} eller en \code{shift}-puls, och visar utgångarna
efter den. Det här är samma sekvens som \code{tx_shift_reg_tb} kontrollerar i sitt första test, så
din egen modul måste återge den rad för rad.

\begin{textdrawing}
          databit  tx_bit  stuff  bit_valid  done  consecutive
ladda           1       1      0          1     0            1
skifta 1        2       1      0          1     0            2
skifta 2        3       1      0          1     0            3
skifta 3        4       1      0          1     0            4
skifta 4        5       1      0          1     0            5
skifta 5        -       0      1          1     0            1
skifta 6        6       0      0          1     0            2
skifta 7        7       0      0          1     0            3
skifta 8        8       0      0          1     0            4
skifta 9        -       0      0          0     1            4
\end{textdrawing}

Tio presentationer för att sända åtta bitar. Var och en av de fyra reglerna syns i den:
\begin{itemize}
  \item \textbf{\code{load} presenterar omedelbart.} Första raden har redan databit 1 på
    \code{tx_bit}, med \code{bit_valid} hög. Ingen \code{shift} har hänt ännu. Hade \code{load} bara
    låst gruppen skulle den här raden fortfarande visa den \emph{föregående} gruppens sista bit, och
    hela ramen skulle gå en bitperiod försent.
  \item \textbf{Följdräkningen driver allt.} Den klättrar från 1 till 5 över de fem \code{'1'}:orna,
    och vid skiftning 4 når den fem. Det lägger inte ut någon stoppbit; det \emph{armerar} en till
    nästa presentation.
  \item \textbf{\code{stuff} är en presentation bred och bär inga data.} Skiftning 5 driver
    \code{not last_bit}, markerar \code{stuff = '1'} och nollställer följden till 1.
    Databitkolumnen visar \code{-}: gruppen har inte gjort några framsteg, vilket är varför åtta
    databitar behöver nio bitläggande presentationer här. \code{bit_valid} är fortfarande \code{'1'},
    eftersom en bit verkligen lades ut på ledningen; den var bara inte en av dina.
  \item \textbf{\code{done} släpar med en, och \code{bit_valid} markerar skillnaden.} Skiftning 8
    lägger ut den sista riktiga biten, och \code{done} förblir låg, eftersom den biten ännu inte fått
    sin bitperiod på bussen. Skiftning 9 är den avslutande skiftningen: \code{tx_bit} håller sitt
    tidigare värde, ingenting nytt läggs ut, \code{bit_valid} faller till \code{'0'}, och först nu
    pulsar \code{done}.
\end{itemize}

Den sista raden är den som är värd att stirra på. \code{tx_bit} läser fortfarande \code{'0'}, precis
som vid skiftning 8, så vad som helst som grindar enbart på \code{shift} skulle mata \code{crc15} med
samma bit en andra gång och tyst korrumpera checksumman. \code{bit_valid} är den enda utgången som
skiljer ”en bit lades ut” från ”en period förflöt”, vilket är varför den finns och varför
\chapref{c16:ch} grindar CRC:n på den.

Lägg slutligen märke till att \code{consecutive} slutar på 4 \textbf{på en följd av \code{'0'}:or},
och \textbf{inte} nollställs. Nästa \code{load} ärver både räkningen och \code{last_bit}, så om den
gruppen öppnar med en \code{'0'} når följden fem redan vid sin allra första bit och en stoppbit
tvingas fram före dess andra. Öppnar den i stället med en \code{'1'} bryts följden och räkningen
börjar om på 1. Det är samma mekanism som den kedjade omladdningen i testbänkens test 3, där en grupp
i stället ärver en följd av fyra \code{'1'}:or och öppnar med en \code{'1'}.

% ------------------------------------------------------------------------------------------
\section{Vad testbänken låser fast, och att lägga in modulen}
\label{c14:sec:tb}

\begin{itemize}
  \item \code{"11111000"} ger en stoppbit (\code{'0'}) direkt efter den femte \code{'1'}:an; en följd
    på exakt fyra stoppas \textbf{inte}.
  \item Att ladda \code{"10101010"} direkt efter en grupp som slutade med fyra \code{'1'}:or (utan
    reset emellan) tvingar fram en stoppbit före den här gruppens andra riktiga bit, vilket bevisar
    att följdtillståndet lever kvar över \code{load}.
  \item \code{bit_valid} är \code{'1'} vid load och vid varje skiftning med riktig bit eller
    stoppbit, \code{'0'} vid den avslutande skiftningen; \code{done} pulsar bara vid just den
    avslutande skiftningen.
  \item \code{bit_valid} är tillbaka på \code{'0'} cykeln \emph{efter} varje presentation, så en
    version som håller nivån underkänns även om den lägger ut varje bit korrekt.
\end{itemize}

\subsection{Att lägga in den i \code{can_controller}}
\label{c14:sub:wire}

Tredje blocket, samma redigering. Signalgruppen i \file{controller/can_controller.vhd} är större här,
eftersom \code{tx_shift_reg} har mer att säga:

\begin{vhdlcode}
    -- tx_shift_reg.
    signal txsr_load                                          : std_logic;
    signal txsr_data                                          : byte_t;
    signal txsr_bit_count                                     : std_logic_vector(3 downto 0);
    signal txsr_shift                                         : std_logic;
    signal txsr_tx_bit, txsr_stuff, txsr_done, txsr_bit_valid : std_logic;
\end{vhdlcode}

\begin{vhdlcode}
    tx_shift_reg1: entity work.tx_shift_reg
        port map(clock, reset_s2_n, txsr_load, txsr_data, txsr_bit_count, txsr_shift,
                 txsr_tx_bit, txsr_stuff, txsr_done, txsr_bit_valid);
\end{vhdlcode}

En enda instans, återanvänd för varje stoppat fält i varje ram. Om det fortfarande ser förvånande ut,
läs om kapitlets inledning: registret laddas om med en ny byte och ett nytt \code{bit_count} per
fält, vilket är precis det som gör att en kopia räcker. (Ta \code{reset_s2_n}-kopplingen som den är
tills vidare; \chapref{c17:ch} byter den mot en reset per ram, av ett skäl bara arbitreringen kan
visa.)
